%\documentclass[envcountsect,11pt,trans]{beamer}
\documentclass[11pt]{article}
\usepackage[english]{babel}
\usepackage[ansinew]{inputenc}
\usepackage{times}
\usepackage{xcolor}
\usepackage{graphicx}
\usepackage{rotating}
\usepackage{bbding,pifont} % two dingbat fonts
\usepackage{amsmath}
\usepackage{amsfonts}
\usepackage{amssymb}
\usepackage{fancybox}
\usepackage{epstopdf}
\usepackage{eso-pic}
\usepackage{ngerman}
\usepackage{tcolorbox}
\usepackage{colortbl}
\usepackage{ulem}
\usepackage{tikz}
\usepackage{booktabs}
\usepackage{longtable}
\usepackage{hyperref}
\usepackage{array}
\usepackage{appendixnumberbeamer}
\usepackage[T1]{fontenc}
\usepackage{adjustbox}
\usepackage{listings}
\usepackage{uarial}
\usepackage{subfigure}
\usepackage{listings}
\usepackage{xcolor}
\usepackage{fancyhdr}
\usetikzlibrary{arrows.meta,positioning,shapes.geometric}

% GIZ color scheme
\definecolor{gizred}{RGB}{204,7,30} % GIZ red
\definecolor{gizgreen}{RGB}{0,104,56} % GIZ green
\definecolor{lightgray}{RGB}{245,245,245}
\definecolor{codegray}{RGB}{248,248,248}

\lstdefinestyle{dynare}{
	basicstyle=\ttfamily\scriptsize,
	backgroundcolor=\color{codegray},
	frame=single,
	rulecolor=\color{gray!40},
	keywordstyle=\color{gizred}\bfseries,
	commentstyle=\color{gray},
	stringstyle=\color{gizgreen},
	columns=fullflexible,
	keepspaces=true,
	showstringspaces=false,
	breaklines=true,
	morekeywords={var,varexo,parameters,model,end,initval,shocks,steady,stoch_simul,stderr,periods,drop,order,irf}
}

\renewcommand{\familydefault}{\sfdefault}

\newcommand*\circled[1]{\tikz[baseline=(char.base)]{
            \node[shape=circle,draw,inner sep=2pt] (char) {#1};}}
\newcommand{\E}{\operatorname{E}}
\newcommand{\Var}{\operatorname{Var}}
\newcommand{\overbar}[1]{\mkern 1.5mu\overline{\mkern-1.5mu#1\mkern-1.5mu}\mkern 1.5mu}

\DeclareMathOperator*{\argmin}{arg\,min}
\DeclareMathOperator*{\argmax}{arg\,max}

\definecolor{Black}{RGB}{0,0,0}
\definecolor{RedGIZ}{RGB}{200,15,15}

% Header and footer for article documents
\pagestyle{fancy}
\fancyhf{}
\chead{\includegraphics[width=2cm,keepaspectratio]{../pictures/GIZ_Logo.png}}
\rhead{\includegraphics[width=2cm,keepaspectratio]{../pictures/IWH_Logo_RGB_DE_Grossformat.png}}
\lfoot{DGE-METRIC Advanced Online Training}
\rfoot{\thepage}
\renewcommand{\headrulewidth}{0.4pt}
\renewcommand{\footrulewidth}{0.4pt}


\title{DGE-METRIC Advanced Online Training}
\author{Christoph Schult} 
\date{June 10--12, 2026}

\begin{document}

\begin{titlepage}
    \centering
    \vspace*{1cm}
    
    % Logos at the top
    \begin{minipage}[t]{3cm}
        \centering
        \includegraphics[width=5cm,keepaspectratio]{../pictures/GIZ_Logo.png}
    \end{minipage}
    \hfill
    \begin{minipage}[t]{3cm}
        \centering
        \includegraphics[width=5cm,keepaspectratio]{../pictures/IWH_Logo_RGB_DE_Grossformat.png}
    \end{minipage}
    \hfill
    {\color{RedGIZ}\Huge \textbf{\inserttitle}}
    
    \vspace{1cm}
    
    {\Large Explicitly AI-Supported Model Development}
    
    \vspace{2cm}
    
    {\large \textbf{Institution}}\\
    Halle Institute for Economic Research (IWH) and GIZ
    
    \vspace{3cm}
    
    \includegraphics[width=cm,keepaspectratio]{../pictures/seperationbar.png}
\end{titlepage}

\newpage

\begin{abstract}
This document outlines a comprehensive advanced online training program for DGE-METRIC users. The training emphasizes explicitly AI-supported model development using modern tools like VS Code and collaborative workflows. The program spans 4--5 days of interactive sessions, totaling 10--15 hours of instruction and hands-on practice.
\end{abstract}

\section*{Training Overview}

\textbf{Institution:} Halle Institute for Economic Research and GIZ

\textbf{Target Audience:} Advanced users with basic experience in Dynare and macroeconomic modeling.

\textbf{Format:} Online, interactive sessions (suggested: 4--5 days, 2--3 hours per session)

\section*{Appetizer: VS Code and AI-Assisted Development}

\subsection*{Why This Matters}

In modern macroeconomic research, efficiency and reproducibility are paramount. This training introduces you to powerful tools that accelerate your workflow and enhance your research productivity.

\subsection*{Key Tools}

\begin{itemize}
    \item \textbf{Visual Studio Code (VS Code):} A lightweight, extensible code editor that streamlines your Dynare/MATLAB work with syntax highlighting, debugging, integrated terminals, and version control integration.
    
    \item \textbf{AI-Assisted Development:} Leverage AI tools (e.g., GitHub Copilot, Claude, and similar assistants) to:
    \begin{itemize}
        \item Generate and refine model equations and code snippets
        \item Automate repetitive coding tasks
        \item Debug errors and optimize performance
        \item Document your model logic and methodology
        \item Accelerate research workflows and reduce development time
    \end{itemize}
\end{itemize}

\subsection*{Your Benefits}

\begin{itemize}
    \item Faster model development and iteration cycles
    \item Better code organization, documentation, and maintainability
    \item Collaborative workflows with version control (Git)
    \item AI-enhanced productivity and research impact
    \item Skills applicable across macroeconomic modeling projects
\end{itemize}

\section*{5-Day Training Agenda}

\subsection*{Day 1: Deep Dive into DGE-METRIC Structure}

\begin{enumerate}
    \item Recap of DGE-METRIC core blocks and equations
    \item Model calibration: advanced parameterization and data sources
    \item Customizing the model for country- or sector-specific analysis
\end{enumerate}

\textbf{Learning Outcomes:} Participants will understand the complete structure of the DGE-METRIC model, master calibration techniques, and be able to adapt the model to country-specific contexts.

\subsection*{Day 2: Advanced Scenario Design}

\begin{enumerate}
    \item Building complex policy scenarios (multi-policy, shocks, time paths)
    \item Using Excel input files: structure, best practices, and troubleshooting
    \item Hands-on: Creating and loading custom scenarios
\end{enumerate}

\textbf{Learning Outcomes:} Participants will design realistic policy scenarios, work efficiently with input data management, and troubleshoot common scenario development issues.

\subsection*{Day 3: Model Extensions and Sensitivity Analysis}

\begin{enumerate}
    \item Adding new sectors, technologies, or climate modules
    \item Sensitivity analysis: parameter sweeps and uncertainty quantification
    \item Automating batch runs and result collection
\end{enumerate}

\textbf{Learning Outcomes:} Participants will extend the model with custom features, conduct rigorous sensitivity analysis, and automate analysis workflows for efficiency.

\subsection*{Day 4: Advanced Output Analysis}

\begin{enumerate}
    \item Interpreting simulation results: macro, sectoral, and distributional impacts
    \item Visualization: advanced plotting in MATLAB/Excel
    \item Exporting results for policy briefs and reports
\end{enumerate}

\textbf{Learning Outcomes:} Participants will extract meaningful insights from model results, create professional visualizations, and communicate findings effectively to policymakers.

\subsection*{Day 5: Open Lab and Q\&A}

\begin{enumerate}
    \item Participants work on their own cases with instructor support
    \item Troubleshooting, tips, and sharing best practices
    \item Wrap-up and feedback
\end{enumerate}

\textbf{Learning Outcomes:} Participants will receive personalized guidance on their research applications and solidify their understanding through practical problem-solving.

\section*{Prerequisites}

\begin{itemize}
    \item Working Dynare and MATLAB installation on your computer
    
    \item Basic knowledge of DGE-METRIC structure (e.g., from the introductory training or documentation)
    
    \item Familiarity with macroeconomic modeling concepts (stocks, flows, equilibrium conditions)
    
    \item Recommended: Basic Git and VS Code experience (or willingness to learn these during the training)
\end{itemize}

\section*{Training Agenda Summary}

The table below provides a quick reference for all five training days:

\vspace{1cm}
\begin{center}
\small
\begin{tabular}{|p{2cm}|p{12cm}|}
    \hline
    \textbf{Day} & \textbf{Topics and Learning Outcomes} \\
    \hline
    \hline
    \textbf{Day 1} & \textbf{Structure}: Recap core blocks, advanced calibration, model customization \\
    \hline
    \textbf{Day 2} & \textbf{Scenarios}: Complex policy scenarios, Excel data management, hands-on practice \\
    \hline
    \textbf{Day 3} & \textbf{Extensions}: Add sectors/technologies, sensitivity analysis, automation \\
    \hline
    \textbf{Day 4} & \textbf{Analysis}: Result interpretation, advanced visualization, report generation \\
    \hline
    \textbf{Day 5} & \textbf{Lab}: Participant projects, troubleshooting, best practices, feedback \\
    \hline
\end{tabular}
\end{center}

\end{document}
